\begin{ExerciseBlock}
\ExerciseHeading{}{Verständnisfragen}{quantum_optics_checks}
Entscheiden Sie, welche Aussagen zutreffen, und begründen Sie Ihre Entscheidung.
\begin{enumerate}
\item Die maximale kinetische Energie der Fotoelektronen hängt beim Einphotonen-Fotoeffekt direkt von der Lichtintensität ab.\TrueFalseAnswer{falsch}
\item Die De-Broglie-Wellenlänge nimmt mit wachsendem Impuls ab.\TrueFalseAnswer{wahr}
\item Laserstrahlung ist immer breitbandig und umfasst viele verschiedene Wellenlängen.\TrueFalseAnswer{falsch}
\item Die Wellenlänge maximaler spektraler Emission eines schwarzen Körpers ist umgekehrt proportional zu seiner Temperatur.\TrueFalseAnswer{wahr}
\item Aus Wellenlänge und Planck-Konstante lässt sich bei bekannter Lichtgeschwindigkeit die Energie eines Photons bestimmen.\TrueFalseAnswer{wahr}
\item Für optische Verstärkung durch stimulierte Emission ist im üblichen Lasermodell eine Besetzungsinversion erforderlich.\TrueFalseAnswer{wahr}
\item Bei einem strahlenden Einphotonen-Übergang entspricht die Photonenenergie näherungsweise der Differenz der Energieniveaus.\TrueFalseAnswer{wahr}
\item Die Lebensdauer eines angeregten Zustands beeinflusst die natürliche spektrale Breite der Emission.\TrueFalseAnswer{wahr}
\item Beim Einphotonen-Fotoeffekt lässt sich die Grenzwellenlänge durch eine höhere Lichtintensität überwinden.\TrueFalseAnswer{falsch}
\item Unabhängig von der Energie der einfallenden Elektronen treten im Röntgenspektrum immer charakteristische Linien auf.\TrueFalseAnswer{falsch}
\item Die Auflösung eines konventionellen Fernfeld-Lichtmikroskops ist durch die Beugung begrenzt.\TrueFalseAnswer{wahr}
\item Im Bohrschen Atommodell sind für Kreisbahnen beliebige Radien möglich.\TrueFalseAnswer{falsch}
\end{enumerate}
\end{ExerciseBlock}
